\subsection{Absolutely regular or absolutely normal algebras}

DEFINITION 1.--- Let \(k\) be a field and \(A\) a \(k\)-algebra. One says that \(A\) is absolutely regular (resp. absolutely normal) \(^{1}\) if the ring \(\mathrm{A}_{(k')}\) is regular (resp. normal) for every radical extension \(k'\) of finite degree over \(k\).

Every \(k\)-algebra that is absolutely regular (resp. absolutely normal) is regular (resp. normal), as is seen by taking \(k' = k\) in Definition 1. Recall that the \(k\)-algebras A for which the ring \(\mathrm{A}_{(k')}\) is reduced for every radical extension \(k'\) of finite degree over \(k\) are the \(k\)-separable algebras (A, V, p. 117, Theorem 2).

Examples.--- 1) If k is perfect, every regular (resp. normal) k-algebra is absolutely regular (resp. absolutely normal).

\begin{enumerate}
\def\labelenumi{\arabic{enumi})}
\setcounter{enumi}{1}
\tightlist
\item
  Let A be an Artinian \(k\)-algebra. If A is normal, it is reduced, hence isomorphic to a finite product of extensions of \(k\) (A, VIII, § 8, no. 1, Proposition 2). The following conditions are equivalent:
\end{enumerate}

A is separable;

A is absolutely regular;

A is absolutely normal.

Indeed, if A is absolutely normal, it is separable; if A is separable, then for every radical extension \(k'\) of finite degree over \(k\), the ring \(\mathrm{A}_{(k')}\) is reduced and Artinian, hence isomorphic to a product of fields and consequently regular, so that \(A\) is absolutely regular.

If moreover the k-algebra A is of finite degree, the preceding conditions are equivalent to saying that A is étale (A, V, p. 34, Theorem 4).

\begin{enumerate}
\def\labelenumi{\arabic{enumi})}
\setcounter{enumi}{2}
\tightlist
\item
  Let A be a local regular \(k\)-algebra. If the extension \(\kappa_{\mathrm{A}}\) of \(k\) is separable, the algebra A is absolutely regular. Indeed, let \(k'\) be a finite extension of \(k\). The A-module \(\mathrm{A}_{(k')}\) is free. Since it is of finite type, every maximal ideal of \(\mathrm{A}_{(k')}\) lies over \(\mathfrak{m}_{\mathrm{A}}\) (V, § 2, no. 1, Proposition 1). The ring \(\kappa_{\mathrm{A}} \otimes_{\mathrm{A}} \mathrm{A}_{(k')}\), isomorphic to \(\kappa_{\mathrm{A}} \otimes_{k} k'\), is regular (no. 3, Lemma 2). By the corollary of Proposition 9 of § 4, no. 5, the ring \(\mathrm{A}_{(k')}\) is regular, which proves that the \(k\)-algebra A is absolutely regular.
\end{enumerate}

PROPOSITION 6. Let k be a field and A a Noetherian k-algebra.

\begin{enumerate}
\def\labelenumi{\alph{enumi})}
\item
  If A is absolutely regular (resp. absolutely normal, resp. separable), then so is \(S^{-1}A\) for every multiplicative subset S of A.
\item
  If \(A_{m}\) is absolutely regular (resp. absolutely normal, resp. separable) for every maximal ideal m of A, then A is absolutely regular (resp. absolutely normal, resp. separable).
\item
  This follows from the fact that the ring \((\mathrm{S}^{-1}\mathrm{A})_{(k^{\prime})}\) is isomorphic to a ring of fractions of the ring \(\mathrm{A}_{(k^{\prime})}\) for every extension \(k^{\prime}\) of k.
\item
  Suppose that \(A_{m}\) is absolutely regular (resp. absolutely normal, resp. separable) for every maximal ideal m of A. Let \(k'\) be a radical extension of k of finite degree and let \(m'\) be a maximal ideal of \(\mathrm{A}_{(k')}\). It remains to verify that the local ring \((\mathrm{A}_{(k')})_{\mathfrak{m}'}\) is regular (resp. normal, resp. reduced).
\end{enumerate}

The canonical homomorphism \(A \to A_{(k')}\) makes \(A_{(k')}\) into a finite A-algebra. The ideal \(m'\) lies over a maximal ideal m of A (V, § 2, no. 1, Proposition 1), and \((\mathrm{A}_{(k')})_{\mathfrak{m}'}\) is isomorphic to a ring of fractions of the regular (resp. normal, resp. reduced) ring \((\mathrm{A}_{\mathfrak{m}})_{(k')}\), hence is regular (resp. normal, resp. reduced).

Lemma 3.--- Let k be a field and K a finitely generated extension of k. There exists an extension L of K, of finite degree, and a subextension k' of L over k, radical and of finite degree over k, such that the extension L of k' is separable.

Choose a composite extension E of K and of a perfect closure \(\hat{k}\) of \(k\), and elements \((t_1, \ldots, t_n)\) of E such that E is a separable algebraic extension of \(\hat{k}(t_1, \ldots, t_n)\). Let I be a finite generating set of K over \(k\). The field \(\hat{k}(t_1, \ldots, t_n)\) (resp. \(\hat{k}\mathrm{K}\)) is the union of the subfields \(k'(t_1, \ldots, t_n)\) (resp. \(k'\mathrm{K}\)), where \(k'\) ranges over the set of subextensions of \(\hat{k}\) of finite degree over \(k\); one can therefore find such a subextension \(k'\) such that the elements \(t_1, \ldots, t_n\) of \(\mathrm{E} = \hat{k}\mathrm{K}\) belong to \(k'\mathrm{K}\) and each element of I is algebraic and separable over \(k'(t_1, \ldots, t_n)\). Then \(\mathrm{L} = k'\mathrm{K}\) is a separable extension of \(k'\) of finite degree over K, and \(k'\) is a radical extension of finite degree over \(k\).

PROPOSITION 7. Let \(k\) be a field, and let A and B be \(k\)-algebras, one of which is essentially of finite type. Suppose that A is absolutely regular (resp. absolutely normal). If the ring B is regular (resp. normal), then so is the ring A \(\otimes_{k}\) B.

Let K be a finitely generated extension of k; let us prove that the ring \(\mathrm{A}_{(\mathrm{K})}\) is regular (resp. normal). Indeed, consider extensions L and \(k'\) having the properties of Lemma 3. Then the ring \(\mathrm{A}_{(k')}\) is regular (resp. normal) by definition, and the ring \(\mathrm{A}_{(\mathrm{L})}\), which is identified with \(\mathrm{A}_{(k')} \otimes_{k'} \mathrm{L}\), is regular (resp. normal) by Proposition 5, b) of No. 3. Consequently, \(\mathrm{A}_{(\mathrm{K})}\) is regular (resp. normal) by Proposition 5, a).

Suppose the ring B is regular (resp. normal) and let us prove the proposition. The canonical homomorphism \(B \to A \otimes_{k} B\) makes \(A \otimes_{k} B\) a free B-module. For every prime ideal p of B, the ring \((A \otimes_{k} B) \otimes_{B} \kappa(\mathfrak{p})\) is identified with \(A \otimes_{k} \kappa(\mathfrak{p})\) by the corollary of Proposition 9 of § 4, No. 5 (resp. Corollary 3 of Theorem 4 of § 1, No. 10); it suffices for us to prove that \(A \otimes_{k} \kappa(\mathfrak{p})\) is regular (resp. normal) for every prime ideal p of B.

If the \(k\)-algebra B is essentially of finite type, the extension \(\kappa(\mathfrak{p})\) of \(k\) is of finite type and the ring A \(\otimes_k \kappa(\mathfrak{p})\) is regular (resp. normal) by what we saw above. Now suppose the \(k\)-algebra A is essentially of finite type; the ring A \(\otimes_k \kappa(\mathfrak{p})\) is Noetherian and is the union of the increasing filtered family of Noetherian subrings \(A \otimes_k K\), where \(K\) runs through the finitely generated subextensions of \(\kappa(\mathfrak{p})\). These latter rings are regular (resp. normal), and one applies Lemma 1 of No. 3.

COROLLARY 1. Let k be a field. The tensor product of two absolutely regular (resp. absolutely normal) k-algebras, one of which is essentially of finite type, is an absolutely regular (resp. absolutely normal) k-algebra.

Let A and B be two k-algebras satisfying the hypotheses of the corollary. Let \(k'\) be a radical extension of k of finite degree. The ring \(\mathrm{B}_{(k')}\) is regular (resp. normal); the ring \(A \otimes_{k} B_{(k')}\) is regular (resp. normal) by prop. 7, as is the ring \((\mathrm{A} \otimes_{k} \mathrm{B}) \otimes_{k} k'\) which is isomorphic to it.

COROLLARY 2.--- Let k be a field, A an absolutely regular (resp. absolutely normal) k-algebra, and K an extension of k. Suppose that A is essentially of finite type or that the extension K of k is of finite type.

\begin{enumerate}
\def\labelenumi{\alph{enumi})}
\item
  The ring \(\mathrm{A}_{(K)}\) is regular (resp. normal).
\item
  If the extension K of k is separable, the k-algebra \(\mathrm{A}_{(\mathrm{K})}\) is absolutely regular (resp. absolutely normal).
\end{enumerate}

Assertion (a) follows from Proposition 7; assertion (b) follows from Corollary 1 and Example 2.

COROLLARY 3. Let \(k\) be a field, let A be a \(k\)-algebra, and let K be an extension of \(k\). Suppose that the \(k\)-algebra A is essentially of finite type or that the extension K of \(k\) is of finite type. For the \(k\)-algebra A to be absolutely regular (resp. absolutely normal), it is necessary and sufficient that the K-algebra \(\mathrm{A}_{(\mathrm{K})}\) be absolutely regular (resp. absolutely normal).

Suppose A is absolutely regular (resp. absolutely normal) and let \(K'\) be a radical extension of K of finite degree. The ring \(K'\otimes_{K}A_{(K)}\), which is isomorphic to \(K'\otimes_{k}A\), is regular (resp. normal) by Corollary 2.

Conversely, suppose the K-algebra \(\mathrm{A}_{(\mathrm{K})}\) is absolutely regular (resp. absolutely normal) and let \(k'\) be a radical extension of \(k\) of finite degree. Let L be a composite extension of \(k'\) and K; then the ring \(\mathrm{A}_{(\mathrm{L})}\) is identified with \(\mathrm{L} \otimes_{\mathrm{K}} \mathrm{A}_{(\mathrm{K})}\), hence is regular (resp. normal); consequently, the ring \(\mathrm{A}_{(k')}\) is regular (resp. normal) by Proposition 5, a) of No. 3.

COROLLARY 4.--- Let k be a field, A a k-algebra essentially of finite type, and K an extension of k which is a perfect field. For A to be absolutely regular (resp. absolutely normal), it is necessary and sufficient that the ring \(\mathrm{A}_{(\mathrm{K})}\) be regular (resp. normal).

This follows from Cor. 3 and Example 1.

